\section{18 aprile una festa della Famiglia Dorotea}

Il 18 aprile quest'anno festeggeremo il nostro beato Don Luca Passi, nel giorno della sua festa liturgica. Spesso il 18 aprile cade o la Settimana Santa o nell'ottava di Pasqua, quindi non è possibile celebrare la liturgia propria, che viene poi posticipata il primo giorno disponibile.

\figurawrap{0.5\textwidth}{don_luca}

Questa volta, sabato 18 aprile, celebreremo la festa del nostro Fondatore in modo speciale alla mattina con la S. Messa propria, nella nostra cappella e alla sera nella S. Messa prefestiva la sua memoria sarà presente attraverso le cooperatrici che rinnoveranno la Promessa.

Come altre volte abbiamo detto, il Beato don Luca, prima delle suore, ha fondato l'Opera di S. Dorotea, nel 1815 e poi nel 1838 le suore per essere l'anima dell'Opera.

Sicuramente quando le suore 175 anni fa sono arrivate a Maserada avranno cercato giovani donne che aiutassero in maniera capillare le fanciulle a vivere in modo sano la loro vita di ogni giorno, attraverso il metodo della \textit{Santa amicizia}.

Questo in fondo era ed è lo stile dell'Opera di Santa Dorotea: farsi vicino a chi ha bisogno e sostenerlo nel cammino, nei vari modi che lo Spirito Santo sa suggerire.

La Parola di Dio che ha fatto nascere in Don Luca il progetto dell'Opera è <<Va, e correggi tuo fratello>>. Un correggere inteso nel senso etimologico \textit{cum-regere}, reggersi insieme, cioè farsi compagni di viaggio. Cogliere il grido inespresso di chi ha bisogno e farsene carico. Questo è quanto è richiesto anche oggi a coloro che aderiscono all'Associazione Ecclesiale approvata dalla Chiesa.

\figurawrap{0.4\textwidth}{logo_dorotee}

A Maserada le cooperatrici sono presenti da molti anni, non hanno cartelli per farsi riconoscere, ma sono presenti come lievito nella pasta all'interno della parrocchia nei vari servizi. Si ritrovano periodicamente per un cammino personale di formazione e crescita spirituale per poter avere sempre quella riserva di lievito necessaria per non perdere la propria identità di Battezzate che annunciano il vangelo con la vita.

La Promessa che rinnovano queste nostre sorelle (speriamo che arrivino anche i fratelli\ldots) in che cosa consiste? Lo scopriremo con il prossimo giornalino.

Buona Pasqua

\firma{Le suore dorotee}
